\documentclass[10pt,a4paper]{amsart}
\usepackage[margin=19mm]{geometry}
\usepackage{amsmath,amssymb,mathtools}
\usepackage[scr=rsfs,cal=euler]{mathalfa}
\usepackage{xfrac}
\usepackage[unicode,hidelinks,bookmarks=false]{hyperref}
\newcommand{\ie}{i.e.\ }
\newcommand{\cf}{cf.\ }
\newcommand{\resp}{respectively\ }
\newcommand{\Subsection}{\subsection}
\providecommand{\nobreakdash}{\nobreak\hbox{-}}
\setlength{\emergencystretch}{2em}
\multlinegap0pt
\usepackage{amssymb}
\newtheorem{proposition}{Proposition}
\title{Connections between colored restricted $b$-ary and ordinary partitions}
\author{Karl Dilcher, Larry Ericksen}
\date{Source: arXiv:2609.03248v1 (CC BY 4.0)}
\begin{document}
\maketitle

\noindent\emph{Source and licence.} Connections between colored restricted $b$-ary and ordinary partitions. Version arXiv:2609.03248v1 (CC BY 4.0), distributed by arXiv under the Creative Commons Attribution 4.0 International licence, http://creativecommons.org/licenses/by/4.0/, as recorded in the arXiv licence metadata for this exact version and shown on the abstract page https://arxiv.org/abs/2609.03248v1. Full licence notice: \texttt{licenses/arxiv-2609.03248-CC-BY-4.0.txt}.\par\smallskip

\noindent\textit{Editorial scope.} Counted unit: the article's proposition prop:6.5 (source0.tex lines 856--866). The article's complete original proof of it is delivered with it (source0.tex lines 868--894, one \texttt{proof} environment of 27 lines). No mathematics is written, repaired or re-derived: the statement and the proof are the article's own lines, in the article's own notation, and no retained line is substituted or re-flowed.  Retained uncounted as the source-local context the proof needs: lines 758--760; lines 774--777.  Nothing was omitted from any retained span, so the package carries no omission record.  No retained line contains a citation, so the wrapper carries no bibliography and no bibliography entry.  No retained line contains a figure environment, an image-inclusion command or drawing code, so no image is delivered and the package declares no asset dependency.  Editorial formatting only, in addition: an AMS \texttt{amsart} preamble that restates the article's own declarations and macros used by the retained lines, namely the environments \texttt{proposition} on separate counters, together with the standard packages those lines need.  The retained environments print in this document as Proposition 1 in order of appearance, whereas the article prints the same items with the numbering of its own full document.  The complete native source of the article as distributed in the arXiv e-print for this exact version is bundled unchanged as \texttt{sources/209281/source0.tex}.
\begin{equation}\label{5.9}
F_b^\rho(q;Z) =\left(\sum_{n=0}^\infty Y_n(Z)q^n\right)F_b^\rho(q^b;Z).
\end{equation}

\begin{equation}\label{6.1}
\sum_{n=0}^\infty Y_nq^n = \left(1+q+q^2+\dots\right)^\rho=\frac{1}{(1-q)^\rho}
=\sum_{n=0}^\infty\binom{n+\rho-1}{\rho-1}q^n,
\end{equation}

\begin{proposition}\label{prop:6.5}
Let $b\geq 2$ and $\rho\geq 1$ be integers. Then the sequence
$\left(C_b^\rho(n)\right)_{n\geq 0}$ of $\rho$-color $b$-ary partitions of $n$
satisfies the identities
\begin{align}
&\sum_{i=0}^\rho(-1)^i\binom{\rho}{i}C_b^\rho(bn-i)=C_b^\rho(n),\label{6.9}\\
&\sum_{i=0}^\rho(-1)^i\binom{\rho}{i}C_b^\rho(bn+r-i)=0\quad (1\leq r\leq b-1),
\label{6.10}
\end{align}
where we assume that $C_b^\rho(m)=0$ for $m<0$.
\end{proposition}

\begin{proof}
Since we have $z_{i,j}=1$ for all $i$ and $j$, we can combine \eqref{6.1} with 
\eqref{5.9} and expand $(1-q)^\rho$, obtaining
\begin{equation}\label{6.11}
\left(\sum_{m=0}^\infty C_b^\rho(m)q^m\right)
\left(\sum_{j=0}^\rho(-1)^j\binom{\rho}{j}q^j\right)
= \sum_{m=0}^\infty C_b^\rho(m)q^{bm}.
\end{equation}
We now take the Cauchy product on the left, setting $n=m+j$. Then \eqref{6.11}
gives
\begin{equation}\label{6.12}
\sum_{n=0}^\infty\left(\sum_{j=0}^\rho(-1)^j\binom{\rho}{j}C_b^\rho(n-j)\right)q^n
= \sum_{m=0}^\infty C_b^\rho(m)q^{bm},
\end{equation}
with the understanding that $C_b^\rho(\ell)=0$ for $\ell< 0$. By equating the
coefficients of equal powers of $q$ in \eqref{6.12}, we get the following two
cases:

(i) When $n=bm$, $m=0, 1, 2,\ldots$, we get
\[
\sum_{j=0}^\rho(-1)^j\binom{\rho}{j}C_b^\rho(bm-j) = C_b^\rho(m),
\]
which is the same as \eqref{6.9}.

(ii) Similarly, when $n=bm+r$, $1\leq r\leq b-1$, the corresponding sum will
be 0, which is the same as \eqref{6.10}.
\end{proof}

\end{document}
